% New English translation from checked cote 25, site 4 (PDF page 5).
% AI-translated by GPT-6 Astra (OpenAI, ChatGPT Pro mode).
\scanpage{25/4}
\mspage{\hand{1}}
IV 20 \quad \uline{Hyperplane sections and \var{\deleted{linear}}{}{deletion}{ega25_25_004-v0001}\var{\inserted{conical}}{conical}{insertion}{ega25_25_004-v0002} projections}.

20.1. \quad \uline{Preliminaries \var{\over{of}{and}}{and}{writing over}{ega25_25_004-v0003} notation}.

Let S be a \quad prescheme, $\underline{E}$ a loc free Module of finite type on S,
$\overset{\hand{\vee}}{\underline{E}}$ its dual. We shall denote by P=$\underline{P}(\underline{E})$ the projective bundle defined by $\underline{E}$,
and by $\overset{\hand{\vee}}{P}$ the one defined by $\overset{\hand{\vee}}{\underline{E}}$. We shall \unsure{call} \ $\overset{\hand{\vee}}{P}$ \ the \uline{scheme of hyperplanes}
of P; this terminology is justified as follows. \ Let
$\hand{\xi}$ be a section of $\overset{\hand{\vee}}{P}$ over S, which is thus\hand{/}determined by an invertible
quotient module L of $\overset{\hand{\vee}}{\underline{E}}$. It gives an invertible quotient module
$\underline{L}_P$ of $(\overset{\hand{\vee}}{\underline{E}})_P$ = $(\underline{E}_P)^{\hand{\vee}}$; on the other hand, we have an invertible quotient module
$\underline{O}_P(1)$ of \ $\underline{E}_P$. Passing \ to the duals, we may regard \ $\underline{L}_P^{-1}$ resp. $\underline{O}_P(-1)$
as invertible submodules (loc direct factors) of \ $\underline{E}_P$ resp.
of $(\underline{E}_P)^{\hand{\vee}}$, \ and the pairing \ $\underline{E}_P \otimes \overset{\hand{\vee}}{\underline{E}}_P \hand{\longrightarrow} \underline{O}_P$ \ therefore defines a
natural pairing
(x) \[ \underline{O}_P(-1)\otimes\underline{L}_P^{-1} \hand{\to} \underline{O}_P \]
or again a transpose \ homomorphism
\var{\margin{$\|$}}{\placedmargin{$\|$}}{marginal note placed}{ega25_25_004-v0004}
(xx) \[ \underline{O}_P \hand{\longrightarrow} \underline{O}_P(1)\otimes\underline{L}_P \ = \ \underline{D}_P(1) \ , \]
i.e. a section of $\underline{L}_P(1)$, canonically determined by $\hand{\xi}$. The ``divisor''
of this section, \ i.e. the \var{\inserted{closed}}{closed}{insertion}{ega25_25_004-v0005} subscheme $H_{\hand{\xi}}$ of \ P defined by the Ideal
that is the image of (x), \ is called the \uline{hyperplane in P} \ defined by \ the element
$\hand{\xi} \hand{\in} \overset{\hand{\vee}}{P}(S)$. \ It can also be described by observing that locally on S \var{\xout{the}}{}{typewriter strike-out}{ega25_25_004-v0006}
\var{\xout{\unsure{datum of $\check{\xi}$ is equivalent}}}{}{typewriter strike-out}{ega25_25_004-v0007} \quad $\hand{\xi}$ \ is given by a section \ $\phi$ of \ $\underline{E}$, \var{\xout{of}}{}{typewriter strike-out}{ega25_25_004-v0008} such
\var{\xout{\illegible{1}}}{}{typewriter strike-out}{ega25_25_004-v0009} that $\phi(s) \neq 0$ for every s \ ($\phi$ being determined by $\hand{\xi}$ \ modulo
multiplication by an invertible section of $\underline{O}_S$); since $\underline{E}$ = $p_{\hand{*}}(\underline{O}_P(1))$
(p: P$\hand{\to}$S being the projection), \ $\phi$ \ can be regarded as a section
of \ $\underline{O}_P(1)$, \var{\over{therefore the}{whose}}{whose}{writing over}{ega25_25_004-v0010} divisor is precisely $H_{\hand{\xi}}$. \ Of course, when we
\var{\inserted[typed]{\uline{regard}}}{\uline{regard}}{insertion}{ega25_25_004-v0011}\var{\xout{\illegible{1}}}{}{typewriter strike-out}{ega25_25_004-v0012} $\underline{L}^{-1}$ \ as an invertible sub\var{\xout{\illegible{1}}}{}{typewriter strike-out}{ega25_25_004-v0013}module of \ $\underline{E}$ that is locally
a direct factor in $\underline{E}$, the correspondence between $\hand{\xi}$ \ (i.e. $\underline{L}$ \hand{\unsure{$\xi$,}}
\tnote{25-004-1}{The abbreviations loc and resp. are retained for locally and respectively. The doubtful n between Nous and appelerons is represented by a doubt mark on call. In (xx) the source has D, while the following sentence has L; both letters are retained. The handwritten c-to-t correction in donc/dont changes the English phrase from therefore the to whose.}

\mspage{\hand{4}}
